% ====================================================================
% Bergin: Mathematics for Economists / 经济学家数学基础
% Chapters 1–3: Introduction, Matrices and Systems of Equations,
%                Linear Algebra: Applications
% Bilingual edition (English + Chinese)
% ====================================================================

% ====================================================================
\chapter{Introduction / 绪论}
% ====================================================================

% -------------------------------------------------------------------
\section{Introduction / 绪论}
% -------------------------------------------------------------------

\begin{original}
This introductory chapter provides an overview of the mathematical
tools that will be developed in the text.  The focus is on those areas
of mathematics most relevant to economic analysis: linear algebra,
multivariate calculus, optimization, and dynamics.
\end{original}
\begin{translation}
本章为绪论，概述本书将要展开的数学工具。重点放在与经济分析最为相关的
数学领域：线性代数、多元微积分、最优化以及动态分析。
\end{translation}

% -------------------------------------------------------------------
\section{A Brief Summary by Chapter / 各章内容简介}
% -------------------------------------------------------------------

\begin{original}
A brief summary by chapter is provided to guide the reader through the
structure of the book.  Each chapter builds on material developed
earlier, so the ordering is important for those reading the book
sequentially.
\end{original}
\begin{translation}
本节提供各章内容简介，帮助读者把握全书的整体结构。各章内容以前面章节
为基础逐步展开，因此对于按顺序阅读的读者而言，章节的编排顺序很重要。
\end{translation}


% ====================================================================
\chapter{Matrices and Systems of Equations / 矩阵与方程组}
% ====================================================================

% -------------------------------------------------------------------
\section{Introduction / 引言}
% -------------------------------------------------------------------

\begin{original}
This chapter introduces matrices and their use in representing and
solving systems of linear equations.  Linear equations and matrix
algebra are fundamental tools in economic modelling, appearing in
everything from simple supply-and-demand models to sophisticated
econometric estimation.
\end{original}
\begin{translation}
本章介绍矩阵及其在表示与求解线性方程组中的应用。线性方程与矩阵代数
是经济建模的基本工具，从简单的供需模型到复杂的计量经济学估计，
到处都有它们的身影。
\end{translation}

% -------------------------------------------------------------------
\subsection{Some Motivating Examples / 若干引例}
% -------------------------------------------------------------------

\begin{original}
We begin with several simple economic examples that illustrate why
matrices and systems of linear equations arise naturally.  These
examples include a simple market equilibrium problem, a national-income
model, and an input–output production model.
\end{original}
\begin{translation}
我们从几个简单的经济学例子入手，说明矩阵与线性方程组何以自然地出现。
这些例子包括一个简单的市场均衡问题、一个国民收入模型以及一个
投入—产出生产模型。
\end{translation}

% -------------------------------------------------------------------
\section{Vectors / 向量}
% -------------------------------------------------------------------

\begin{original}
Vectors are ordered lists of numbers.  Geometrically, a vector can be
viewed as a point or a directed line segment in $n$-dimensional space.
Algebraically, vectors are the building blocks of matrix operations.
\end{original}
\begin{translation}
向量是有序数组。从几何上看，向量可视为 $n$ 维空间中的一个点或一条
有向线段。从代数上看，向量是矩阵运算的基本构件。
\end{translation}

% -------------------------------------------------------------------
\section{Matrices / 矩阵}
% -------------------------------------------------------------------

\begin{original}
A matrix is a rectangular array of numbers.  The size or order of a
matrix is given by the number of rows and columns it contains.  An
$m \times n$ matrix has $m$ rows and $n$ columns.
\end{original}
\begin{translation}
矩阵是一个矩形的数表。矩阵的大小或阶由其行数和列数决定。
一个 $m \times n$ 矩阵有 $m$ 行和 $n$ 列。
\end{translation}

% -------------------------------------------------------------------
\subsection{Matrix Addition and Subtraction / 矩阵加法与减法}
% -------------------------------------------------------------------

\begin{original}
Two matrices of the same dimension can be added (or subtracted) by
adding (or subtracting) corresponding entries.  If $A$ and $B$ are both
$m \times n$, then $C = A + B$ is also $m \times n$, with
$c_{ij} = a_{ij} + b_{ij}$.
\end{original}
\begin{translation}
两个维数相同的矩阵可以通过对应元素相加（或相减）来进行加法
（或减法）运算。若 $A$ 与 $B$ 均为 $m \times n$ 矩阵，则
$C = A + B$ 也是 $m \times n$ 矩阵，且 $c_{ij} = a_{ij} + b_{ij}$。
\end{translation}

% -------------------------------------------------------------------
\subsection{Matrix Multiplication / 矩阵乘法}
% -------------------------------------------------------------------

\begin{original}
The product $C = AB$ is defined when the number of columns of $A$
equals the number of rows of $B$.  If $A$ is $m \times r$ and $B$ is
$r \times n$, then $C$ is $m \times n$ with entries
$c_{ij} = \sum_{k=1}^{r} a_{ik} b_{kj}$.
\end{original}
\begin{translation}
当 $A$ 的列数等于 $B$ 的行数时，乘积 $C = AB$ 有定义。若 $A$ 为
$m \times r$ 矩阵，$B$ 为 $r \times n$ 矩阵，则 $C$ 为 $m \times n$
矩阵，其元素为 $c_{ij} = \sum_{k=1}^{r} a_{ik} b_{kj}$。
\end{translation}

% -------------------------------------------------------------------
\subsection{Matrix Transpose / 矩阵转置}
% -------------------------------------------------------------------

\begin{original}
The transpose of an $m \times n$ matrix $A$, denoted $A'$ or $A^{T}$,
is the $n \times m$ matrix obtained by interchanging rows and columns:
$(A')_{ij} = A_{ji}$.
\end{original}
\begin{translation}
$m \times n$ 矩阵 $A$ 的转置，记作 $A'$ 或 $A^{T}$，是一个
$n \times m$ 矩阵，通过互换行与列得到：$(A')_{ij} = A_{ji}$。
\end{translation}

% -------------------------------------------------------------------
\subsection{Matrix Inversion / 矩阵求逆}
% -------------------------------------------------------------------

\begin{original}
A square matrix $A$ is invertible (or nonsingular) if there exists a
matrix $A^{-1}$ such that $AA^{-1} = A^{-1}A = I$, where $I$ is the
identity matrix.  Not every square matrix has an inverse; those that
do not are called singular.
\end{original}
\begin{translation}
若存在矩阵 $A^{-1}$ 使得 $AA^{-1} = A^{-1}A = I$（其中 $I$ 为单位矩阵），
则方阵 $A$ 可逆（或称非奇异）。并非每个方阵都有逆矩阵；没有逆矩阵的
方阵称为奇异矩阵。
\end{translation}

% -------------------------------------------------------------------
\subsection{Additional Remarks on Determinants / 关于行列式的补充说明}
% -------------------------------------------------------------------

\begin{original}
The determinant of a square matrix, denoted $\det(A)$ or $|A|$, is a
scalar that provides information about the matrix.  A square matrix is
invertible if and only if its determinant is nonzero.  The determinant
also has a geometric interpretation as the volume scaling factor of the
linear transformation represented by the matrix.
\end{original}
\begin{translation}
方阵的行列式，记作 $\det(A)$ 或 $|A|$，是一个能够提供矩阵信息的标量。
方阵可逆当且仅当其行列式非零。行列式还具有几何解释，即矩阵所表示的
线性变换的体积缩放因子。
\end{translation}

% -------------------------------------------------------------------
\subsection{Cramer's Rule / 克莱姆法则}
% -------------------------------------------------------------------

\begin{original}
Cramer's rule gives an explicit formula for the solution of a system of
$n$ linear equations in $n$ unknowns, provided the coefficient matrix
has nonzero determinant.  The $i$-th variable is $x_i = |A_i|/|A|$,
where $A_i$ is the matrix obtained by replacing the $i$-th column of
$A$ with the right-hand-side vector $b$.
\end{original}
\begin{translation}
克莱姆法则给出了 $n$ 个方程 $n$ 个未知数的线性方程组解的显式公式，
前提是系数矩阵的行列式非零。第 $i$ 个变量为 $x_i = |A_i|/|A|$，
其中 $A_i$ 是将 $A$ 的第 $i$ 列替换为右端向量 $b$ 后得到的矩阵。
\end{translation}

% -------------------------------------------------------------------
\subsection{Gaussian Elimination / 高斯消元法}
% -------------------------------------------------------------------

\begin{original}
Gaussian elimination is a systematic procedure for solving systems of
linear equations by reducing the augmented matrix to row-echelon form
(or reduced row-echelon form) through a sequence of elementary row
operations.
\end{original}
\begin{translation}
高斯消元法是一种求解线性方程组的系统化方法，通过一系列初等行变换
将增广矩阵化为行阶梯形（或简化行阶梯形）。
\end{translation}

% -------------------------------------------------------------------
\section{Linear Dependence and Independence / 线性相关与线性无关}
% -------------------------------------------------------------------

\begin{original}
A set of vectors $\{v_1, v_2, \ldots, v_k\}$ is linearly dependent if
there exist scalars $c_1, c_2, \ldots, c_k$, not all zero, such that
$c_1 v_1 + c_2 v_2 + \cdots + c_k v_k = 0$.  Otherwise the vectors are
linearly independent.
\end{original}
\begin{translation}
向量组 $\{v_1, v_2, \ldots, v_k\}$ 称为线性相关的，若存在不全为零的
标量 $c_1, c_2, \ldots, c_k$ 使得
$c_1 v_1 + c_2 v_2 + \cdots + c_k v_k = 0$。
否则，这些向量是线性无关的。
\end{translation}

% -------------------------------------------------------------------
\subsection{Matrix Rank and Linear Independence / 矩阵的秩与线性无关性}
% -------------------------------------------------------------------

\begin{original}
The rank of a matrix is the maximum number of linearly independent
rows (or columns) in the matrix.  Row rank and column rank are always
equal.  The rank provides information about the dimension of the
solution space of the associated homogeneous system.
\end{original}
\begin{translation}
矩阵的秩是矩阵中线性无关行（或列）的最大数目。行秩与列秩始终相等。
秩提供了关于相应齐次方程组解空间维数的信息。
\end{translation}

% -------------------------------------------------------------------
\section{Solutions of Systems of Equations / 方程组的解}
% -------------------------------------------------------------------

\begin{original}
A system of linear equations may have no solution, exactly one
solution, or infinitely many solutions.  The existence and uniqueness
of solutions depend on the relationship between the rank of the
coefficient matrix and the rank of the augmented matrix.
\end{original}
\begin{translation}
一个线性方程组可能无解，恰有一解，或有无穷多解。解的存在性与唯一性
取决于系数矩阵的秩与增广矩阵的秩之间的关系。
\end{translation}

% -------------------------------------------------------------------
\subsection{Solving Systems of Equations: A Geometric View / 方程组求解：几何视角}
% -------------------------------------------------------------------

\begin{original}
Geometrically, a system of linear equations can be viewed as the
intersection of hyperplanes.  Each equation defines a hyperplane, and
the solution set is the intersection of all these hyperplanes.  This
geometric perspective helps in understanding the possible solution
configurations.
\end{original}
\begin{translation}
从几何上看，线性方程组可视为超平面的交。每个方程定义一个超平面，
解集就是所有这些超平面的交集。这一几何视角有助于理解可能的解的
各种情形。
\end{translation}

% -------------------------------------------------------------------
\subsection{Solutions of Equation Systems: Rank and Linear Independence / 方程组的解：秩与线性无关性}
% -------------------------------------------------------------------

\begin{original}
The solution structure of $Ax = b$ is determined by the rank of $A$ and
the rank of the augmented matrix $[A \mid b]$.  When these ranks are
equal, solutions exist; the difference between the number of unknowns
and the rank gives the dimension of the solution space for the
homogeneous system.
\end{original}
\begin{translation}
$Ax = b$ 的解的结构由 $A$ 的秩和增广矩阵 $[A \mid b]$ 的秩决定。
当这两个秩相等时，解存在；未知数个数与秩的差给出了齐次方程组解空间
的维数。
\end{translation}

% -------------------------------------------------------------------
\section{Special Matrices / 特殊矩阵}
% -------------------------------------------------------------------

\begin{original}
Several classes of matrices have special properties that are useful in
economic applications: symmetric matrices ($A' = A$), idempotent
matrices ($A^2 = A$), orthogonal matrices ($A'A = I$), diagonal
matrices, triangular matrices, and positive definite matrices.
\end{original}
\begin{translation}
有几类矩阵因其特殊性质在经济应用中非常有用：对称矩阵 ($A' = A$)、
幂等矩阵 ($A^2 = A$)、正交矩阵 ($A'A = I$)、对角矩阵、三角矩阵
以及正定矩阵。
\end{translation}

% -------------------------------------------------------------------
\section*{Exercises for Chapter 2 / 第 2 章习题}
% -------------------------------------------------------------------

\begin{original}
Exercises for Chapter 2 cover matrix operations, determinants, solving
systems of equations using Gaussian elimination and Cramer's rule,
linear dependence, rank, and special matrices.
\end{original}
\begin{translation}
第 2 章习题涵盖矩阵运算、行列式、运用高斯消元法和克莱姆法则求解
方程组、线性相关性、秩以及特殊矩阵等内容。
\end{translation}


% ====================================================================
\chapter{Linear Algebra: Applications / 线性代数：应用}
% ====================================================================

% -------------------------------------------------------------------
\section{Introduction / 引言}
% -------------------------------------------------------------------

\begin{original}
This chapter applies the linear-algebra techniques developed in
Chapter~2 to several economic problems.  The applications range from
market models and production analysis to dynamic models of inflation
and econometric estimation.
\end{original}
\begin{translation}
本章将第 2 章展开的线性代数工具应用于若干经济问题。应用范围涵盖市场
模型和生产分析，以及通货膨胀动态模型与计量经济学估计。
\end{translation}

% -------------------------------------------------------------------
\section{The Market Model / 市场模型}
% -------------------------------------------------------------------

\begin{original}
A simple linear market model expresses supply and demand as linear
functions of price.  The equilibrium price vector solves a system of
linear equations.  Matrix algebra provides a compact and powerful
framework for analysing multi-market equilibrium.
\end{original}
\begin{translation}
简单的线性市场模型将供给和需求表示为价格的线性函数。均衡价格向量是
一组线性方程组的解。矩阵代数为分析多市场均衡提供了一个紧凑而强大的
框架。
\end{translation}

% -------------------------------------------------------------------
\section{Adjacency Matrices / 邻接矩阵}
% -------------------------------------------------------------------

\begin{original}
An adjacency matrix represents the connections in a network or graph.
In economics, adjacency matrices appear in the analysis of social
networks, trade linkages, and industrial organisation.  Powers of the
adjacency matrix give information about indirect connections.
\end{original}
\begin{translation}
邻接矩阵表示网络或图中的连接关系。在经济学中，邻接矩阵出现在社会网络、
贸易关联和产业组织分析中。邻接矩阵的幂提供了关于间接连接的信息。
\end{translation}

% -------------------------------------------------------------------
\section{Input--Output Analysis / 投入—产出分析}
% -------------------------------------------------------------------

\begin{original}
Input–output analysis, developed by Wassily Leontief, describes the
interdependence among the different sectors of an economy.  The
input–output matrix records the flows of goods between sectors.  The
Leontief inverse $(I - A)^{-1}$ is central to computing the total
(direct and indirect) requirements needed to meet a given final demand.
\end{original}
\begin{translation}
投入—产出分析由 Wassily Leontief 提出，描述了经济中不同部门之间的
相互依赖关系。投入—产出矩阵记录了部门间商品的流动。Leontief 逆矩阵
$(I - A)^{-1}$ 对于计算满足给定最终需求所需的总（直接和间接）需求
是核心的。
\end{translation}

% -------------------------------------------------------------------
\section{Inflation and Unemployment Dynamics / 通货膨胀与失业动态}
% -------------------------------------------------------------------

\begin{original}
Dynamic models of inflation and unemployment can be written as systems
of linear difference equations.  The long-run behaviour of the system
depends on the eigenvalues of the coefficient matrix.  This section
shows how matrix methods illuminate the dynamics of macroeconomic
variables.
\end{original}
\begin{translation}
通货膨胀与失业的动态模型可以写成线性差分方程组的形式。系统的长期行为
取决于系数矩阵的特征值。本节展示矩阵方法如何揭示宏观经济变量的动态
特征。
\end{translation}

% -------------------------------------------------------------------
\section{Stationary (Invariant) Distributions / 平稳（不变）分布}
% -------------------------------------------------------------------

\begin{original}
A stationary (or invariant) distribution of a Markov chain is a
probability distribution $\pi$ satisfying $\pi P = \pi$, where $P$ is
the transition matrix.  Such distributions describe the long-run
behaviour of the chain and are important in many economic applications,
including income mobility and firm-size dynamics.
\end{original}
\begin{translation}
马尔可夫链的平稳（或不变）分布是一个概率分布 $\pi$，满足
$\pi P = \pi$，其中 $P$ 为转移矩阵。这类分布描述了链的长期行为，
在许多经济应用中都很重要，包括收入流动性和企业规模动态。
\end{translation}

% -------------------------------------------------------------------
\subsection{Convergence to a Stationary Distribution / 收敛到平稳分布}
% -------------------------------------------------------------------

\begin{original}
Under appropriate conditions (irreducibility and aperiodicity), a
Markov chain converges to its unique stationary distribution regardless
of the initial distribution.  The speed of convergence is governed by
the second-largest eigenvalue of the transition matrix.
\end{original}
\begin{translation}
在适当的条件（不可约性和非周期性）下，马尔可夫链无论从何种初始分布
出发，都会收敛到其唯一的平稳分布。收敛速度由转移矩阵的第二大特征值
决定。
\end{translation}

% -------------------------------------------------------------------
\subsection{Computation of the Invariant Distribution / 不变分布的计算}
% -------------------------------------------------------------------

\begin{original}
The invariant distribution can be computed by solving the linear system
$\pi (I - P) = 0$ subject to the normalisation constraint
$\sum_i \pi_i = 1$.  Standard matrix methods, including Gaussian
elimination, can be used to find $\pi$.
\end{original}
\begin{translation}
不变分布可以通过求解线性方程组 $\pi (I - P) = 0$ 并附加归一化约束
$\sum_i \pi_i = 1$ 来计算。标准矩阵方法，包括高斯消元法，可用于
求出 $\pi$。
\end{translation}

% -------------------------------------------------------------------
\section{Econometrics / 计量经济学}
% -------------------------------------------------------------------

\begin{original}
Linear algebra is the mathematical foundation of econometrics.  The
linear regression model $y = X\beta + \varepsilon$ is expressed in
matrix notation, and the least-squares estimator is derived using
matrix differentiation and projection.
\end{original}
\begin{translation}
线性代数是计量经济学的数学基础。线性回归模型 $y = X\beta + \varepsilon$
用矩阵记号表示，最小二乘估计量则通过矩阵微分与投影推导得出。
\end{translation}

% -------------------------------------------------------------------
\subsection{Derivation of the Least Squares Estimator / 最小二乘估计量的推导}
% -------------------------------------------------------------------

\begin{original}
The ordinary least squares (OLS) estimator minimises the sum of squared
residuals $e'e = (y - X\beta)'(y - X\beta)$.  The first-order
condition yields the normal equations $X'X\hat{\beta} = X'y$, and,
provided $X'X$ is invertible, the OLS estimator is
$\hat{\beta} = (X'X)^{-1}X'y$.
\end{original}
\begin{translation}
普通最小二乘 (OLS) 估计量极小化残差平方和
$e'e = (y - X\beta)'(y - X\beta)$。一阶条件给出正规方程
$X'X\hat{\beta} = X'y$，且只要 $X'X$ 可逆，OLS 估计量即为
$\hat{\beta} = (X'X)^{-1}X'y$。
\end{translation}

% -------------------------------------------------------------------
\section*{Exercises for Chapter 3 / 第 3 章习题}
% -------------------------------------------------------------------

\begin{original}
Exercises for Chapter 3 cover the market model, adjacency matrices,
input–output analysis, inflation and unemployment dynamics, stationary
distributions, and the least-squares estimator.
\end{original}
\begin{translation}
第 3 章习题涵盖市场模型、邻接矩阵、投入—产出分析、通货膨胀与失业
动态、平稳分布以及最小二乘估计量等内容。
\end{translation}
